\documentclass[11pt]{article}
\usepackage[T1]{fontenc}
\usepackage{lmodern}
\usepackage[margin=1in]{geometry}
\usepackage{amsmath,amssymb,amsthm,mathtools}
\usepackage{microtype}
\usepackage[hidelinks]{hyperref}
\usepackage[nameinlink,capitalise,noabbrev]{cleveref}
\usepackage{authblk}
\renewcommand{\Affilfont}{\small}

\newtheorem{theorem}{Theorem}[section]
\newtheorem{proposition}[theorem]{Proposition}
\newtheorem{lemma}[theorem]{Lemma}
\newtheorem{corollary}[theorem]{Corollary}
\theoremstyle{remark}
\newtheorem{remark}[theorem]{Remark}

% The environments above share the theorem counter, so cleveref would name
% them all "Theorem"; these aliases restore the correct names.
\usepackage{etoolbox}
\AtBeginEnvironment{proposition}{\crefalias{theorem}{proposition}}
\AtBeginEnvironment{lemma}{\crefalias{theorem}{lemma}}
\AtBeginEnvironment{corollary}{\crefalias{theorem}{corollary}}
\AtBeginEnvironment{remark}{\crefalias{theorem}{remark}}

\newcommand{\ind}{\mathbf{1}}
\newcommand{\E}{\mathbb{E}}
\newcommand{\R}{\mathbb{R}}
\newcommand{\dsq}{\delta_{\square}}
\newcommand{\Tr}{\operatorname{Tr}}
\newcommand{\op}{\mathrm{op}}
\newcommand{\norm}[1]{\left\lVert #1\right\rVert}
\newcommand{\abs}[1]{\left\lvert #1\right\rvert}
\newcommand{\Var}{\operatorname{Var}}

\title{Goodman-type lower bounds for odd cycles III: Quantitative stability}
\author[1,2]{Taeyoung~Kim}
\author[3]{Jineon~Baek}
\author[4]{Seonghyuk~Im}
\author[2]{Joonkyung~Lee}
\author[1,2]{Hongseok~Yang}
\affil[1]{School of Computing, KAIST, Korea}
\affil[2]{School of Computational Sciences, Korea Institute for Advanced Study (KIAS), Korea}
\affil[3]{June E Huh Center for Mathematical Challenges, Korea Institute for Advanced Study (KIAS), Korea}
\affil[4]{Center for AI and Natural Science, Korea Institute for Advanced Study (KIAS), Korea}
\affil[ ]{\newline\texttt{taeyoungkim21@kaist.ac.kr}, \texttt{{jineon.kias}@gmail.com}, \texttt{seonghyuk@kias.re.kr}, \texttt{joonkyunglee@kias.re.kr}, \texttt{hongseokyang@kias.re.kr}}
\date{September 23, 2026}

\begin{document}
\maketitle

\begin{abstract}
Assuming the results proved before the stability corollary in
\emph{Goodman-type lower bounds for odd cycles I: The dense regime},
we prove the following. Fix an odd integer $m\geq3$ and an integer
$k\geq3$, and let $p_k=1-1/k$. If a graphon $W$ has edge density
within $\delta_e$ of $p_k$ and $C_m$-density at most
$\delta_C$ above the extremal value, then $W$ lies within
$O_{m,k}\bigl(\sqrt{\delta_e+\delta_C}\bigr)$ of the balanced complete
$k$-partite graphon, in relabelled $L^1$ distance and hence in cut
distance. The exponent $1/2$ is optimal in each of the two parameters
separately: it is already forced by graphons at exact edge density,
and independently by graphons whose $C_m$-density does not exceed the
extremal value at all. No constant is tracked; every implied constant
depends only on $m$ and $k$.
\end{abstract}

\section{Statement}
\label{sec:statement}

References to numbered results of the manuscript mean the numbering in
\cite{kim2026goodman}. The argument below is conditional on its Theorem~1.1 and
Theorem~1.3, its excursion expansion, its polynomial nonnegativity, and
its Gamma-moment lemma; it uses neither its Proposition~6.1 nor its
qualitative stability corollary.

All graphons are symmetric measurable functions
$[0,1]^2\to[0,1]$, and integrals with unspecified domains are over the
appropriate powers of $[0,1]$. For a graphon $W$ we write
\[
 p=t(K_2,W),\qquad U=1-W,\qquad q=1-p,
 \qquad d_W(x)=\int_0^1W(x,y)\,dy,
\]
and $(T_Wf)(x)=\int_0^1W(x,y)f(y)\,dy$ for the associated integral
operator. The path $P_j$ has $j$ edges. Throughout, $m\geq3$ is odd and
\[
 G_m(p)=p^m-p(1-p)^{m-1}.
\]
For a bounded measurable kernel $F$ we use the cut norm and cut
distance~\cite{lovasz2006limits,borgs2008convergent,lovasz2012large},
$\norm{F}_{\square}=\sup_{S,T}\abs{\int_{S\times T}F}$ over measurable
$S,T\subseteq[0,1]$, and
\[
 \dsq(W,V)=\inf_\phi\norm{W-V^\phi}_{\square},
 \qquad
 \delta_1(W,V)=\inf_\phi\norm{W-V^\phi}_1,
\]
where $V^\phi(x,y)=V(\phi(x),\phi(y))$ and $\phi$ ranges over
measure-preserving bijections of $[0,1]$ modulo null sets. Directly from
these definitions,
\begin{equation}\label{eq:distance-basics}
 \dsq(W,V)\leq\delta_1(W,V)\leq1,
 \qquad
 \delta_1(1-W,1-V)=\delta_1(W,V),
\end{equation}
and both quantities satisfy the triangle inequality, by composing the
relabellings and using the invariance of the two norms under them.

For an integer $k\geq3$ we set $q_k=1/k$ and $p_k=1-1/k$, we let $W_k$ be
the balanced complete $k$-partite graphon, and we put $U_k=1-W_k$.

\paragraph{Background.}
Stability statements of this kind go back to Simonovits~\cite{simonovits1968method}. For clique densities, the structure of graphs that nearly minimize the density was determined by Pikhurko and Razborov~\cite{pikhurko2017asymptotic} for triangles and by Kim, Liu, Pikhurko and Sharifzadeh~\cite{kim2020asymptotic} for general cliques. For odd cycles, \cite[Corollary~6.3]{kim2026goodman} gives qualitative stability at the densities $p_k$; the present article makes it quantitative.

\paragraph{Convention on constants.}
Throughout, $m\geq3$ is a fixed odd integer and $k\geq3$ a fixed integer.
We write $A=O(B)$, or $A\lesssim B$, to mean $A\leq CB$ for some
$C=C(m,k)<\infty$ depending only on $m$ and $k$; the value of $C$ may
change from line to line. No constant is tracked.

\begin{theorem}[Two-parameter stability]
\label{thm:main}
Fix an odd integer $m\geq3$ and an integer $k\geq3$. Let
$\delta_e,\delta_C\geq0$ and let $W$ be a graphon with
\begin{equation}\label{eq:hypotheses}
 \abs{t(K_2,W)-p_k}\leq\delta_e,
 \qquad
 t(C_m,W)\leq G_m(p_k)+\delta_C.
\end{equation}
Then
\begin{equation}\label{eq:main}
 \dsq(W,W_k)\ \leq\ \delta_1(W,W_k)
 \ \lesssim\ (\delta_e+\delta_C)+\sqrt{\delta_e+\delta_C}
 \ \lesssim\ \sqrt{\delta_e+\delta_C}
 \ \lesssim\ \sqrt{\delta_e}+\sqrt{\delta_C}.
\end{equation}
\end{theorem}

The middle bound in \eqref{eq:main} is the informative one: for small
tolerances the square-root term dominates, and its coefficient is
considerably smaller than the coefficient of the single-power bound,
which must also cover the range where the tolerances are of order one.

\begin{theorem}[Optimality in each parameter separately]
\label{thm:sharp}
Fix an odd integer $m\geq3$ and an integer $k\geq3$, and let
$c>1/2$. There is no constant $C<\infty$ for which
\[
 \dsq(W,W_k)\leq C(\delta_e+\delta_C)^c
\]
holds for all sufficiently small $\delta_e,\delta_C\geq0$ and all
graphons $W$ satisfying \eqref{eq:hypotheses}. This already fails on each
of the two coordinate axes:
\begin{enumerate}
\item[(i)] there are graphons with $t(K_2,W)=p_k$ exactly, hence with
$\delta_e=0$, admissible $\delta_C\downarrow0$, and
$\dsq(W,W_k)\gtrsim\sqrt{\delta_C}$;
\item[(ii)] there are graphons with $t(C_m,W)\leq G_m(p_k)$, hence with
$\delta_C=0$, admissible $\delta_e\downarrow0$, and
$\dsq(W,W_k)\gtrsim\sqrt{\delta_e}$.
\end{enumerate}
\end{theorem}

In particular neither parameter may be given an exponent larger than
$1/2$, so $\sqrt{\delta_e+\delta_C}$ is the correct joint form. The two
families witnessing (i) and (ii) are the same complete $k$-partite
graphons with one enlarged and one reduced part; only the normalisation
differs.

The proof of \Cref{thm:main} occupies
Sections~\ref{sec:degree}--\ref{sec:completion} and has three
ingredients, each stated as a self-contained estimate:

\begin{itemize}
\item a \emph{degree estimate} (\Cref{prop:degree}): the gap in the
two-colour inequality of \cite[Theorem~1.3]{kim2026goodman} dominates the
variance of the degree function, with a positive constant depending only
on $m$;
\item \emph{robust spectral estimates} (\Cref{lem:spectral}): a cycle
count close to the extremal value forces the complement to be nearly
$\{0,1\}$-valued and nearly transitive, quantitatively;
\item a \emph{rounding lemma} (\Cref{lem:rounding}): a $\{0,1\}$-valued
graphon with nearly correct degrees and few cherries with a missing
closing edge is close to a balanced union of $k$ cliques.
\end{itemize}

The third step is purely combinatorial and constructive. Its selection
step is the pivot step of the correlation-clustering algorithm of Ailon,
Charikar and Newman~\cite{ailon2008aggregating}, and it uses neither the
regularity lemma~\cite{szemeredi1978regular} nor the induced removal
lemma~\cite{alon2000efficient}; its only smallness threshold
is explicit, with the trivial bound $\delta_1\leq1$ covering the
complementary range. \Cref{thm:sharp} is proved in
Section~\ref{sec:sharpness}.

\section{The degree estimate}
\label{sec:degree}

Recall from \cite[Section~3]{kim2026goodman} that for $d\geq0$, $h_d$ denotes the
complete homogeneous symmetric polynomial of degree $d$, with $h_0=1$,
and that for $p=1-q$ and $n=m-2$ the first excursion polynomial is
\begin{equation}\label{eq:first-poly}
 P_{m,1}(q;\lambda)
 =m\bigl[h_{n}(p,-\lambda)+h_{n}(q,\lambda)\bigr]
   -h_{n-1}(q,q,\lambda).
\end{equation}

\begin{lemma}[Strict positivity of the first excursion polynomial]
\label{lem:first-poly}
Let
\[
 c_m=\min\bigl\{P_{m,1}(q;\lambda):\ 0\leq q\leq\tfrac13,\
 \abs{\lambda}\leq\tfrac12\bigr\}.
\]
Then $c_m>0$.
\end{lemma}

\begin{proof}
Write $n=m-2$, which is odd. Expanding \eqref{eq:first-poly} by
$h_{n}(p,-\lambda)=\sum_{j=0}^{n}(-1)^jp^{n-j}\lambda^j$,
$h_{n}(q,\lambda)=\sum_{j=0}^{n}q^{n-j}\lambda^j$ and
$h_{n-1}(q,q,\lambda)=\sum_{j=0}^{n-1}(n-j)q^{n-1-j}\lambda^j$ gives
\begin{equation}\label{eq:poly-expansion}
 P_{m,1}(q;\lambda)
 =\sum_{j=0}^{n-1}
 \Bigl[m\bigl((-1)^jp^{n-j}+q^{n-j}\bigr)
       -(n-j)q^{n-1-j}\Bigr]\lambda^j,
\end{equation}
the degree-$n$ terms cancelling because $n$ is odd. For every odd $j$
the bracket equals
$m\bigl(q^{n-j}-p^{n-j}\bigr)-(n-j)q^{n-1-j}$, which is nonpositive
because $p\geq q\geq0$. Consequently
\begin{equation}\label{eq:reflection}
 P_{m,1}(q;-a)-P_{m,1}(q;a)
 =-2\sum_{j\ \mathrm{odd}}
 \Bigl[m\bigl(q^{n-j}-p^{n-j}\bigr)-(n-j)q^{n-1-j}\Bigr]a^j\ \geq\ 0
 \qquad(a\geq0),
\end{equation}
so it suffices to bound $P_{m,1}(q;\lambda)$ from below for
$\lambda\geq0$.

Let $\lambda>0$. With $\rho_{n,m}$ as in \cite[Eq.~(5.7)]{kim2026goodman}, the
Gamma representation \cite[Eq.~(5.15)]{kim2026goodman}, specialised to $r=1$,
reads
\[
 P_{m,1}(q;\lambda)
 =n\,\E\Bigl[\rho_{n,m}\Bigl(q+\frac{\lambda}{Z}Y\Bigr)\Bigr],
 \qquad Y\sim\operatorname{Gamma}(1,1),\quad
 Z\sim\operatorname{Gamma}(m-1,1)
\]
with $Y$ and $Z$ independent. The centred identity
\cite[Eq.~(5.18)]{kim2026goodman}, again with $r=1$, expresses
$n\rho_{n,m}(\tfrac12+v)$ as $2a_0$ plus a combination, with positive
coefficients $a_j=2^{2j+1-n}\binom{n}{2j}$, of the quantities
$2(1+j)v^{2j}-jv^{2j-1}$ for $1\leq j\leq (n-1)/2$; here
$a_0=2^{1-n}$. For $0\leq q\leq1/3$ and any $x\geq0$, the proof of
\cite[Proposition~5.6]{kim2026goodman} shows that each of these quantities has
nonnegative expectation at $v=q+xY-\tfrac12$. Hence
\[
 n\,\E\rho_{n,m}(q+xY)\ \geq\ 2a_0=2^{2-n}
 \qquad(x\geq0,\ 0\leq q\leq\tfrac13),
\]
and conditioning on $Z$ gives $P_{m,1}(q;\lambda)\geq2^{2-n}>0$ for every
$\lambda>0$. Continuity extends this to $\lambda=0$, and
\eqref{eq:reflection} extends it to $\lambda<0$. Thus
$c_m\geq2^{2-n}=2^{4-m}>0$.
\end{proof}

Following \cite[Eq.~(6.1)]{kim2026goodman}, we write, for an arbitrary graphon
$W$ of edge density $p$ with $U=1-W$ and $q=1-p$,
\begin{equation}\label{eq:Phi-definition}
 \Phi_m(W)=t(C_m,W)+t(C_m,U)-t(P_{m-1},U)-G_m(p),
\end{equation}
so that \cite[Theorem~1.3]{kim2026goodman}, a strengthening of the
extended-commonality inequality of Kim and Lee~\cite{kim2024extended} for
odd cycles, states $\Phi_m(W)\geq0$ for $p\geq2/3$, and \cite[Eq.~(6.3)]{kim2026goodman} is the gap decomposition
\begin{equation}\label{eq:gap-decomposition}
 t(C_m,W)-G_m(p)
 =\Phi_m(W)+\bigl[t(P_{m-1},U)-t(C_m,U)\bigr],
\end{equation}
in which, for $p\geq2/3$, both terms on the right are nonnegative.

\begin{proposition}[Degree estimate]\label{prop:degree}
Every graphon $W$ with edge density $p\geq2/3$ satisfies
\begin{equation}\label{eq:degree-coercivity}
 \Phi_m(W)\ \geq\ c_m\int_0^1\bigl(d_W(x)-p\bigr)^2\,dx .
\end{equation}
\end{proposition}

\begin{proof}
Suppose first that $W$ is a step graphon as in
\cite[Section~2]{kim2026goodman}. The excursion expansion
\cite[Eq.~(3.10)]{kim2026goodman} reads
\[
 \Phi_m(W)=\sum_{r=1}^{(m-1)/2}
 \int P_{m,r}(q;\lambda_1,\ldots,\lambda_r)
       \,d\mu(\lambda_1)\cdots d\mu(\lambda_r),
\]
where the measure $\mu$ is nonnegative and, by
\cite[Lemma~2.2 and Eq.~(2.3)]{kim2026goodman}, is supported on $[-1/2,1/2]$ and
has total mass
\[
 \mu(\R)=\norm{g}_2^2=\int_0^1\bigl(d_W(x)-p\bigr)^2\,dx .
\]
Since $p\geq2/3$ gives $q\leq1/3$, all the integrands are nonnegative by
\cite[Section~5]{kim2026goodman}. Discarding every term with $r\geq2$ and applying
\Cref{lem:first-poly} to the term $r=1$ yields
\eqref{eq:degree-coercivity}.

For a general graphon let $W_s$ be the same-density dyadic step
approximations of \cite[Lemma~2.1 and Appendix~A]{kim2026goodman}, so that
$\norm{W_s-W}_1\to0$ while $t(K_2,W_s)=p$ for every $s$. Telescoping the
edge factors shows that every homomorphism density occurring in
\eqref{eq:Phi-definition} converges, and the degree variance converges as
well because it equals $t(P_2,W_s)-p^2$. Passing to the limit proves the
assertion.
\end{proof}

\begin{remark}
\Cref{prop:degree} is exactly the stability statement for
\cite[Theorem~1.3]{kim2026goodman}: the gap in the two-colour inequality is
bounded below by a fixed multiple of the degree variance, and it vanishes
precisely for $p$-regular graphons. Retaining instead only the term
$r=(m-1)/2$, whose integrand is the constant $m-1$, recovers
\cite[Proposition~6.1]{kim2026goodman}. Retaining every term gives the full
lower bound $\Phi_m(W)\geq\sum_r c_{m,r}\sigma^{2r}$ with
$c_{m,r}=\min P_{m,r}$ over $[-1/2,1/2]^r$, together with the matching
upper bound obtained by replacing each minimum by the corresponding
maximum; we shall not need either refinement.
\end{remark}

\section{Robust spectral estimates}
\label{sec:spectral}

For a graphon $U$ define
\begin{align}
 \eta(U)&=\int U(x,y)\bigl(1-U(x,y)\bigr)\,dx\,dy,\label{eq:eta-def}\\
 \tau(U)&=t(P_2,U)-t(C_3,U)
 =\int U(x,y)U(y,z)\bigl(1-U(x,z)\bigr)\,dx\,dy\,dz.\label{eq:tau-def}
\end{align}
Both are nonnegative. The first measures how far $U$ is from being
$\{0,1\}$-valued; the second measures how far the relation
$\{U=1\}$ is from being transitive, and vanishes exactly when $U$ is a
disjoint union of cliques.

\begin{lemma}[Even-path inequality]\label{lem:path}
If $U$ has edge density $q$ and $\ell\geq2$ is even, then
$t(P_\ell,U)\geq q^\ell$.
\end{lemma}

This is the even case of the inequality of Mulholland and
Smith~\cite{mulholland1959inequality} and Blakley and
Roy~\cite{blakley1965holder}, which holds for every $\ell$. For even $\ell$
the spectral proof is short.

\begin{proof}
Let $\nu$ be the spectral measure of the self-adjoint operator $T_U$ at
the unit vector $\ind$. Then $\nu$ is a probability measure with
$\int\lambda\,d\nu=\langle\ind,T_U\ind\rangle=q$ and
$\int\lambda^\ell\,d\nu=\langle\ind,T_U^\ell\ind\rangle=t(P_\ell,U)$. As
$\ell$ is even, $x\mapsto x^\ell$ is convex on $\R$, and Jensen's
inequality gives the claim.
\end{proof}

\begin{lemma}[Robust moment comparison]\label{lem:spectral}
Let $U$ be a graphon of edge density $q=1/k$, put
$\sigma^2=\int_0^1(d_U-q)^2$, and suppose that $\Gamma\geq0$ satisfies
\begin{equation}\label{eq:moment-hypothesis}
 t(C_m,U)\geq q^{m-1}-\Gamma .
\end{equation}
Then
\begin{equation}\label{eq:operator-bound}
 \norm{T_U}_{\op}\leq q+\frac{\sigma}{q},
\end{equation}
and
\begin{equation}\label{eq:eta-tau-bound}
 \eta(U)\ \lesssim\ \Gamma+\sigma,
 \qquad
 \tau(U)\ \lesssim\ \Gamma+\sigma .
\end{equation}
\end{lemma}

\begin{proof}
Write $a=m-2$, a positive odd integer, and
$\varrho=\norm{T_U}_{\op}$. The operator $T_U$ is self-adjoint and
Hilbert--Schmidt, hence compact with real spectrum, and
\[
 q=\langle\ind,T_U\ind\rangle\leq\varrho\leq\norm{U}_2\leq1 .
\]

\emph{A nonnegative extremal eigenfunction.}
The following is the argument behind the Perron--Frobenius theorem for
nonnegative matrices~\cite[Chapter~8]{horn2012matrix}. Let $h$ be a real unit eigenfunction for an eigenvalue of modulus
$\varrho$. Since $U\geq0$,
\[
 \varrho=\abs{\langle h,T_Uh\rangle}\leq\langle\abs h,T_U\abs h\rangle
 \leq\varrho,
\]
so $f=\abs h$ is a nonnegative unit vector with
$\langle f,T_Uf\rangle=\varrho$. Then
$\norm{T_Uf-\varrho f}_2^2=\norm{T_Uf}_2^2-\varrho^2\leq0$, whence
$T_Uf=\varrho f$.

\emph{The operator norm.}
From $\varrho f=T_Uf$, $U\leq1$ and $f\geq0$ we get
$\norm f_\infty\leq\varrho^{-1}\norm f_1\leq\varrho^{-1}$ and therefore
$\norm{f^2}_2\leq\norm f_\infty\leq\varrho^{-1}$. Using symmetry and
$2f(x)f(y)\leq f(x)^2+f(y)^2$,
\[
 \varrho=\int U(x,y)f(x)f(y)\,dx\,dy
 \leq\int d_U(x)f(x)^2\,dx
 \leq q+\sigma\norm{f^2}_2\leq q+\frac{\sigma}{\varrho}
 \leq q+\frac{\sigma}{q},
\]
which is \eqref{eq:operator-bound}; in particular
\begin{equation}\label{eq:rho-minus-q}
 0\leq\varrho-q\leq\frac{\sigma}{q}\ \lesssim\ \sigma .
\end{equation}

\emph{The deficit.}
Let $(\lambda_i)$ list the nonzero eigenvalues of $T_U$ with
multiplicity, so that, by the Hilbert--Schmidt identity and the trace
formula~\cite{lovasz2012large}, $\sum_i\lambda_i^2=\int U^2\leq q$ and
$t(C_m,U)=\sum_i\lambda_i^m$; all series below converge absolutely.
Since $\eta(U)=q-\int U^2$,
\begin{equation}\label{eq:F}
 F:=\varrho^aq-t(C_m,U)
 =\varrho^a\eta(U)+\sum_i\lambda_i^2\bigl(\varrho^a-\lambda_i^a\bigr),
\end{equation}
and every summand is nonnegative because $a$ is odd and
$\lambda_i\leq\varrho$. By \eqref{eq:moment-hypothesis},
\begin{equation}\label{eq:F-upper}
 0\leq F\leq\Gamma+q\bigl(\varrho^a-q^a\bigr).
\end{equation}
For $0\leq x\leq1$ one has $1-x^a\leq a(1-x)$, so
$q(\varrho^a-q^a)=q\varrho^a\bigl(1-(q/\varrho)^a\bigr)
\leq a\varrho^{a}(\varrho-q)\lesssim\sigma$ by \eqref{eq:rho-minus-q}
and $\varrho\leq1$. Combining with \eqref{eq:F} and $\varrho\geq q$,
\[
 \eta(U)\leq\varrho^{-a}F\leq q^{-a}\bigl(\Gamma+q(\varrho^a-q^a)\bigr)
 \ \lesssim\ \Gamma+\sigma ,
\]
which is the first half of \eqref{eq:eta-tau-bound}.

\emph{Transitivity.}
For $-1\leq x\leq1$ and odd $a$ one has $1-x\leq2(1-x^a)$: for $x\geq0$
because $x^a\leq x$, and for $x\leq0$ because the left side is at most
$2$ while $1-x^a\geq1$. Applying this with $x=\lambda_i/\varrho$ termwise
in
\[
 \varrho q-t(C_3,U)
 =\varrho\eta(U)+\sum_i\lambda_i^2(\varrho-\lambda_i)
\]
gives $\varrho q-t(C_3,U)\leq2\varrho^{1-a}F$. Since
$t(P_2,U)=q^2+\sigma^2$ and $\varrho\geq q$,
\[
 \tau(U)=\sigma^2+\bigl(q^2-\varrho q\bigr)
 +\bigl(\varrho q-t(C_3,U)\bigr)
 \leq\sigma^2+2q^{1-a}\bigl(\Gamma+q(\varrho^a-q^a)\bigr)
 \ \lesssim\ \sigma^2+\Gamma+\sigma ,
\]
and $\sigma\leq1$ gives $\sigma^2\leq\sigma$. This is the second half of
\eqref{eq:eta-tau-bound}.
\end{proof}

\section{Recovering a balanced clique partition}
\label{sec:rounding}

\begin{lemma}[Approximation by a balanced union of cliques]
\label{lem:rounding}
Let $k\geq2$, let $q=1/k$, and let $H$ be a $\{0,1\}$-valued graphon.
Put
\[
 v=\int_0^1\abs{d_H-q},\qquad \tau=\tau(H).
\]
Then $\delta_1(H,U_k)\lesssim v+\tau$.
\end{lemma}

\begin{proof}
Set $b=v+12k^2\tau$. If $b\geq1/(16k^2)$ then
$16k^2b\geq1\geq\delta_1(H,U_k)$ and there is nothing to prove, so
assume
\begin{equation}\label{eq:small-b}
 b<\frac{1}{16k^2}.
\end{equation}
We show $\delta_1(H,U_k)\leq(16k+3)b$, which suffices. Replacing $H$ on a
null set by a symmetric Borel $\{0,1\}$-valued representative, all sets
below may be taken Borel and no integral is affected.

\paragraph{The cost of selecting a neighbourhood.}
For a measurable $R\subseteq[0,1]$ and $x\in R$ put
$N_R(x)=\{y\in R:H(x,y)=1\}$ and $d_R(x)=\abs{N_R(x)}$, where $\abs A$
denotes Lebesgue measure. Define
\begin{equation}\label{eq:local-cost}
 e_R(x)=\int_{N_R(x)^2}(1-H)
       +2\int_{N_R(x)\times(R\setminus N_R(x))}H,
\end{equation}
the $L^1$ cost of making $N_R(x)$ a clique and of removing its edges to
the rest of $R$. All costs are computed against the original $H$, which
is never modified. Writing the membership indicator of $N_R(x)$ as
$H(x,\cdot)$ and using Fubini's theorem and the symmetry of $H$,
\begin{equation}\label{eq:average-cost}
 \int_Re_R(x)\,dx
 =\int_{R^3}H(x,y)H(x,z)\bigl(1-H(y,z)\bigr)
 +2\int_{R^3}H(x,y)\bigl(1-H(x,z)\bigr)H(y,z)
 =3\tau_R\leq3\tau,
\end{equation}
where $\tau_R$ is the integral \eqref{eq:tau-def} restricted to $R^3$.
Each cherry with a missing closing edge is counted three times, once from
each of its vertices. This is the charging argument for the pivot
algorithm of Ailon, Charikar and Newman~\cite{ailon2008aggregating}, in
which a vertex and its neighbourhood are taken as one cluster.

\paragraph{Extraction.}
Suppose disjoint sets $A_1,\ldots,A_t$ have been selected, let
$R=[0,1]\setminus\bigcup_{i\leq t}A_i$, and let $E$ be the sum of their
costs \eqref{eq:local-cost} at the times of selection. Each edge of $H$
from $R$ to a selected set was counted, together with its symmetric
counterpart, in the cost of that selection, and every other contribution
to $E$ is nonnegative; hence $\int_{R\times R^c}H\leq E/2$. Since
$d_H(x)-d_R(x)=\int_{R^c}H(x,y)\,dy$ for $x\in R$,
\begin{equation}\label{eq:remaining-degree}
 \int_R\abs{d_R-q}\leq v+\int_{R\times R^c}H\leq v+E/2 .
\end{equation}

We select every set to have measure at least $q/2$ and cost at most
$12k\tau$. Suppose inductively that $t$ such choices have been made.
Their disjointness forces $t\leq2k$, so
\begin{equation}\label{eq:total-cost-bound}
 E\leq12kt\tau\leq24k^2\tau,\qquad v+E/2\leq b .
\end{equation}
If $\abs R\geq q/2$, set $X_R=\{x\in R:d_R(x)\geq q/2\}$. On
$R\setminus X_R$ we have $\abs{d_R-q}\geq q/2$, so by
\eqref{eq:remaining-degree} and \eqref{eq:small-b},
$\abs{R\setminus X_R}\leq2k(v+E/2)\leq2kb<q/8$, whence
$\abs{X_R}\geq q/4$. By \eqref{eq:average-cost} there is $x\in X_R$ with
$e_R(x)\leq3\tau/\abs{X_R}\leq12k\tau$. Select $A_{t+1}=N_R(x)$, which
has measure $d_R(x)\geq q/2$, and replace $R$ by $R\setminus A_{t+1}$.
Both inductive properties are preserved, and the induction starts with no
selected sets and $E=0$.

Each choice removes measure at least $q/2$, so the process stops after
at most $2k$ choices. Write the final sets as $A_1,\ldots,A_N$, put
$a_i=\abs{A_i}$, and let $R$ be the final remainder, of measure $r$. Then
\begin{equation}\label{eq:extraction-output}
 a_i\geq q/2,\qquad N\leq2k,\qquad r<q/2,\qquad E\leq24k^2\tau=2(b-v).
\end{equation}

\paragraph{The remainder and the number of parts.}
For $x\in R$ we have $d_R(x)\leq r<q/2$, so
\eqref{eq:remaining-degree} gives $qr/2\leq b$ and hence $r\leq2kb$;
moreover $r^2\leq(q/2)r\leq b$. Put
$V_0=\sum_{i=1}^N\ind_{A_i\times A_i}$. The costs counted during the
extraction account exactly for the discrepancies between $H$ and $V_0$
outside $R^2$, so
\begin{equation}\label{eq:partial-cliques}
 \norm{H-V_0}_1=E+\int_{R^2}H\leq E+r^2\leq3b .
\end{equation}
The degree function of $V_0$ equals $a_i$ on $A_i$ and $0$ on $R$, and
$\norm{d_H-d_{V_0}}_1\leq\norm{H-V_0}_1$, so by
\eqref{eq:extraction-output},
\begin{equation}\label{eq:part-degrees}
 \sum_{i=1}^Na_i\abs{a_i-q}+qr=\int\abs{d_{V_0}-q}
 \leq v+E+r^2\leq3b .
\end{equation}
Since $a_i\geq1/(2k)$, this gives $\sum_i\abs{a_i-q}\leq6kb$, and since
$\sum_ia_i=1-r$,
\[
 \abs{1-Nq}\leq r+\sum_i\abs{a_i-q}\leq8kb<q/2 .
\]
The number $1-Nq=(k-N)/k$ is an integer multiple of $q$, so
\begin{equation}\label{eq:number-of-parts}
 N=k .
\end{equation}

\paragraph{Balancing.}
By atomlessness choose measurable $A_i'\subseteq A_i$ with
$\abs{A_i'}=\min(a_i,q)$. The pool of unretained points has measure
\[
 s=r+\sum_{i=1}^k(a_i-q)_+=\sum_{i=1}^k(q-a_i)_+
 \leq r+\sum_{i=1}^k\abs{a_i-q}\leq8kb,
\]
the middle equality following from $\sum_ia_i=1-r$ and $kq=1$. Split the
pool into sets of measures $(q-a_i)_+$ and adjoin the $i$th to $A_i'$;
this produces a measurable partition $B_1,\ldots,B_k$ of $[0,1]$ into
sets of measure $q$. Put $V_B=\sum_i\ind_{B_i\times B_i}$. On pairs both
of whose endpoints were retained, $V_B$ and $V_0$ agree, so their
difference is supported on pairs meeting the pool and
\[
 \norm{V_B-V_0}_1\leq2s\leq16kb .
\]
Finally $V_B$ is a relabelling of $U_k$: on $B_i$ the map
$\phi(x)=(i-1)q+\int_0^x\ind_{B_i}$ sends Lebesgue measure restricted to
$B_i$ to Lebesgue measure on the $i$th interval of length $q$, is
strictly increasing at the density-one points of $B_i$, and therefore
has a measurable inverse after discarding null sets; combining the $k$
maps gives a measure-preserving bijection modulo null sets with
$U_k^\phi=V_B$ almost everywhere. With \eqref{eq:partial-cliques},
\[
 \delta_1(H,U_k)\leq\norm{H-V_B}_1\leq(16k+3)b
 =(16k+3)\bigl(v+12k^2\tau\bigr)\ \lesssim\ v+\tau. \qedhere
\]
\end{proof}

\section{Proof of the two-parameter bound}
\label{sec:completion}

We first treat graphons whose edge density is exactly $p_k$.

\begin{proposition}[Exact-density estimate]\label{prop:exact}
Let $W$ be a graphon with $t(K_2,W)=p_k$ and set
$\Gamma=t(C_m,W)-G_m(p_k)$. Then $\Gamma\geq0$ and
\begin{equation}\label{eq:exact-estimate}
 \delta_1(W,W_k)\ \lesssim\ \Gamma+\sqrt{\Gamma}.
\end{equation}
\end{proposition}

\begin{proof}
Put $p=p_k$, $q=1/k$ and $U=1-W$, and let
$\sigma^2=\int_0^1(d_W-p)^2$; since $d_U-q=-(d_W-p)$, this also equals
$\int_0^1(d_U-q)^2$. As $p\geq2/3$, \cite[Theorem~1.1]{kim2026goodman} gives
$\Gamma\geq0$.

\emph{Two small quantities.}
By the gap decomposition \eqref{eq:gap-decomposition} both summands are
nonnegative and their sum is $\Gamma$, so
\[
 \Phi_m(W)\leq\Gamma
 \qquad\text{and}\qquad
 0\leq t(P_{m-1},U)-t(C_m,U)\leq\Gamma .
\]
\Cref{prop:degree} turns the first into
\begin{equation}\label{eq:sigma-bound}
 \sigma^2\leq\Gamma/c_m,
 \qquad\text{that is,}\qquad \sigma\lesssim\sqrt{\Gamma},
\end{equation}
and \Cref{lem:path}, applied with the even length $m-1$, turns the
second into $t(C_m,U)\geq t(P_{m-1},U)-\Gamma\geq q^{m-1}-\Gamma$. Hence
\Cref{lem:spectral} applies to $U$ with this $\Gamma$ and this $\sigma$,
and yields
\begin{equation}\label{eq:eta-tau-applied}
 \eta(U)\lesssim\Gamma+\sigma,
 \qquad
 \tau(U)\lesssim\Gamma+\sigma .
\end{equation}

\emph{Thresholding.}
Let $H=\ind_{\{U\geq1/2\}}$ and $e=\norm{H-U}_1$. Pointwise
$\abs{H-U}=\min(U,1-U)\leq2U(1-U)$, so by \eqref{eq:eta-tau-applied},
\begin{equation}\label{eq:threshold-error}
 e\leq2\eta(U)\ \lesssim\ \Gamma+\sigma .
\end{equation}
Since $\norm{d_H-d_U}_1\leq e$ and $\norm{d_U-q}_1\leq\sigma$,
\begin{equation}\label{eq:threshold-degrees}
 v_H:=\int_0^1\abs{d_H-q}\leq\sigma+e\ \lesssim\ \Gamma+\sigma .
\end{equation}
Telescoping the three factors in \eqref{eq:tau-def}, each difference
being bounded by $\abs{H-U}$ on its own pair of coordinates and the
remaining coordinate contributing a factor one, gives
$\abs{\tau(H)-\tau(U)}\leq3e$, so by \eqref{eq:eta-tau-applied} and
\eqref{eq:threshold-error},
\begin{equation}\label{eq:threshold-triples}
 \tau(H)\ \lesssim\ \Gamma+\sigma .
\end{equation}

\emph{Rounding.}
By \eqref{eq:distance-basics}, the triangle inequality, and
\Cref{lem:rounding} applied to the $\{0,1\}$-valued graphon $H$,
\[
 \delta_1(W,W_k)=\delta_1(U,U_k)
 \leq e+\delta_1(H,U_k)
 \ \lesssim\ e+v_H+\tau(H)
 \ \lesssim\ \Gamma+\sigma .
\]
Substituting \eqref{eq:sigma-bound} proves \eqref{eq:exact-estimate}.
\end{proof}

\begin{proof}[Proof of \Cref{thm:main}]
Let $W$ satisfy \eqref{eq:hypotheses} and put $p_0=t(K_2,W)$. Define
\begin{equation}\label{eq:density-correction}
 \widehat W=
 \begin{cases}
 \dfrac{p_k}{p_0}\,W, & p_0\geq p_k,\\[8pt]
 W+\dfrac{p_k-p_0}{1-p_0}\,(1-W), & p_0<p_k .
 \end{cases}
\end{equation}
In either case the denominator is positive and the coefficient lies in
$[0,1]$, so $\widehat W$ is a graphon, and direct integration gives
\begin{equation}\label{eq:density-correction-error}
 t(K_2,\widehat W)=p_k,
 \qquad
 \norm{\widehat W-W}_1=\abs{p_0-p_k}\leq\delta_e .
\end{equation}

For graphons $F_1,F_2$ and a finite simple graph $F$, telescoping the
edge factors in the defining integral gives
$\abs{t(F,F_1)-t(F,F_2)}\leq\abs{E(F)}\norm{F_1-F_2}_1$. Applying this to
$C_m$ and using \eqref{eq:hypotheses},
\begin{equation}\label{eq:Gamma-hat}
 0\ \leq\ \widehat\Gamma:=t(C_m,\widehat W)-G_m(p_k)
 \ \leq\ \delta_C+m\,\delta_e
 \ \lesssim\ \delta_e+\delta_C ,
\end{equation}
the left inequality being \cite[Theorem~1.1]{kim2026goodman} applied to
$\widehat W$, whose edge density is exactly $p_k\geq2/3$. Only
$\widehat W$ is ever subjected to the dense-regime theorem and to
\Cref{prop:degree}; in particular the argument remains valid when $k=3$
and $p_0<2/3$.

Writing $D=\delta_e+\delta_C$ and combining
\eqref{eq:density-correction-error}, \Cref{prop:exact} and
\eqref{eq:Gamma-hat},
\begin{equation}\label{eq:two-term}
 \delta_1(W,W_k)\leq\norm{W-\widehat W}_1+\delta_1(\widehat W,W_k)
 \ \lesssim\ \delta_e+\widehat\Gamma+\sqrt{\widehat\Gamma}
 \ \lesssim\ D+\sqrt D ,
\end{equation}
which is the first bound in \eqref{eq:main}. For the second, recall
$\delta_1(W,W_k)\leq1$ and note that $\min\{1,x\}\leq\sqrt x$ and
$\min\{1,x+y\}\leq\min\{1,x\}+y$ for $x,y\geq0$; hence, with
$x\lesssim D$ and $y\lesssim\sqrt D$ the two terms of
\eqref{eq:two-term},
\[
 \delta_1(W,W_k)\leq\min\bigl\{1,\ C D+C\sqrt D\bigr\}
 \leq\sqrt{CD}+C\sqrt D\ \lesssim\ \sqrt D .
\]
The final bound in \eqref{eq:main} follows from
$\sqrt{x+y}\leq\sqrt x+\sqrt y$, and $\dsq\leq\delta_1$ is
\eqref{eq:distance-basics}.
\end{proof}

\section{Optimality in each parameter}
\label{sec:sharpness}

Fix $m$ and $k$ and write $q=1/k$, $p=1-q$. For $0<\varepsilon<q$ let
$W_\varepsilon$ be the complete $k$-partite graphon with part measures
\[
 q+\varepsilon,\quad q-\varepsilon,\quad q,\ \ldots,\ q .
\]
Its edge density is
\begin{equation}\label{eq:unbalanced-edge}
 t(K_2,W_\varepsilon)
 =1-\bigl((q+\varepsilon)^2+(q-\varepsilon)^2+(k-2)q^2\bigr)
 =p-2\varepsilon^2 .
\end{equation}

\begin{lemma}[Cycle density under a change of part sizes]
\label{lem:expansion}
There are $\beta=\beta(m,k)>0$ and $\varepsilon_0=\varepsilon_0(m,k)>0$
such that
\begin{equation}\label{eq:unbalanced-cycle}
 t(C_m,W_\varepsilon)\leq G_m(p)-\tfrac12\beta\varepsilon^2
 \qquad(0<\varepsilon\leq\varepsilon_0).
\end{equation}
\end{lemma}

\begin{proof}
Let $J$ be the $k\times k$ all-one matrix, $I$ the identity, and put
$D=\operatorname{diag}(1,-1,0,\ldots,0)$, $B_0=q(J-I)$ and $R=(J-I)D$.
Summing over the successive part indices of a cycle gives
$t(C_m,W_{(s_1,\ldots,s_k)})
=\Tr\bigl(((J-I)\operatorname{diag}(s_1,\ldots,s_k))^m\bigr)$, so
\[
 t(C_m,W_\varepsilon)=\Tr\bigl((B_0+\varepsilon R)^m\bigr).
\]
This polynomial in $\varepsilon$ is even, because interchanging the first
two part indices replaces $\varepsilon$ by $-\varepsilon$, and its
constant term is $G_m(p)$ by \cite[Eq.~(1.1)]{kim2026goodman}.

Its coefficient of $\varepsilon^2$ is
$\tfrac m2\sum_{j=0}^{m-2}\Tr\bigl(RB_0^jRB_0^{m-2-j}\bigr)$: there are
$m$ choices for the position of the first factor $\varepsilon R$ and
$j\in\{0,\ldots,m-2\}$ intervening factors $B_0$ before the second,
cyclicity of the trace makes the value independent of the first position,
and each unordered pair of positions is counted twice. Let
$P_0=J/k$ and $P_1=I-P_0$ be the two complementary orthogonal
projections. Then $B_0=pP_0-qP_1$ and $J-I=(k-1)P_0-P_1$, and from
$P_0DP_0=0$, $\Tr(P_0D^2)=2q$ and $\Tr(D^2)=2$ one computes
\[
 P_0RP_0=0,\qquad
 \Tr(P_0RP_1R)=-(k-1)\Tr(P_0DP_1D)=-2p,\qquad
 \Tr(P_1RP_1R)=2(1-2q).
\]
Writing $a=m-2$ and expanding $B_0^j=p^jP_0+(-q)^jP_1$ therefore gives
\[
 \Tr\bigl(RB_0^jRB_0^{a-j}\bigr)
 =-2p\bigl(p^j(-q)^{a-j}+(-q)^jp^{a-j}\bigr)+2(1-2q)(-q)^a .
\]
Since $a$ is odd, $\sum_{j=0}^ap^j(-q)^{a-j}
=\bigl(p^{a+1}-(-q)^{a+1}\bigr)/(p+q)=p^{m-1}-q^{m-1}$ and
$(-q)^a=-q^a$, so the coefficient of $\varepsilon^2$ equals $-\beta$ with
\[
 \beta=2mp\bigl(p^{m-1}-q^{m-1}\bigr)+m(m-1)(1-2q)q^{m-2}>0,
\]
positivity holding because $p>q$ and, as $k\geq3$, $1-2q>0$.

For the higher coefficients, expand the cycle density as the sum over
proper assignments of part indices to the cycle vertices of the product
of the corresponding part measures. A vertex contributing a factor
$\varepsilon$ must receive one of two part indices and contributes a
factor of absolute value one, while each remaining vertex has $k$ choices
and contributes $q$; dropping the properness restrictions therefore
bounds the coefficient of $\varepsilon^j$ in absolute value by
$\binom mj2^jk^{m-j}q^{m-j}=\binom mj2^j$. Odd coefficients vanish, so for
$\varepsilon\leq1$ the terms of degree at least four contribute at most
$\varepsilon^4\sum_{j\geq4}\binom mj2^j\leq3^m\varepsilon^4$ in absolute
value. Taking
$\varepsilon_0=\min\{q/2,\sqrt{\beta/(2\cdot3^m)}\}$ makes this at most
$\tfrac12\beta\varepsilon^2$, which proves
\eqref{eq:unbalanced-cycle}.
\end{proof}

\begin{lemma}[Distance to the balanced graphon]\label{lem:distance}
For $0<\varepsilon<q$ we have
$\dsq(W_\varepsilon,W_k)\geq(q+\varepsilon)\varepsilon$.
\end{lemma}

\begin{proof}
Every relabelling $W_k^\phi$ has constant degree function $p$. Let $S$ be
the part of $W_\varepsilon$ of measure $q+\varepsilon$, on which
$d_{W_\varepsilon}=1-(q+\varepsilon)=p-\varepsilon$. Testing the cut norm
on $S\times[0,1]$ gives, for every $\phi$,
\[
 \norm{W_\varepsilon-W_k^\phi}_{\square}
 \geq\Bigl\lvert\int_S\bigl(d_{W_\varepsilon}-p\bigr)\Bigr\rvert
 =(q+\varepsilon)\varepsilon .
\]
Taking the infimum over $\phi$ proves the claim.
\end{proof}

\begin{proof}[Proof of \Cref{thm:sharp}]
Let $0<\varepsilon\leq\varepsilon_0$ throughout.

\emph{(ii) Pure edge-density tolerance.}
By \Cref{lem:expansion}, $t(C_m,W_\varepsilon)\leq G_m(p_k)$, so
$W_\varepsilon$ satisfies \eqref{eq:hypotheses} with
$\delta_C=0$ and, by \eqref{eq:unbalanced-edge},
$\delta_e=2\varepsilon^2$. By \Cref{lem:distance},
\[
 \dsq(W_\varepsilon,W_k)\geq q\varepsilon
 =\frac{q}{\sqrt2}\sqrt{\delta_e}.
\]
Thus a bound $C(\delta_e+\delta_C)^c$ would force
$q2^{-c}\varepsilon^{1-2c}\leq C$ for all small $\varepsilon$, which
fails for every $c>1/2$.

\emph{(i) Pure cycle tolerance.}
Let
\[
 \widetilde W_\varepsilon
 =W_\varepsilon+\frac{2\varepsilon^2}{q+2\varepsilon^2}
 \bigl(1-W_\varepsilon\bigr).
\]
Since $\int(1-W_\varepsilon)=q+2\varepsilon^2$ by
\eqref{eq:unbalanced-edge}, we get
$t(K_2,\widetilde W_\varepsilon)=p_k$ exactly and
$\norm{\widetilde W_\varepsilon-W_\varepsilon}_1=2\varepsilon^2$.
Telescoping the edge factors of $C_m$ and using \Cref{lem:expansion},
\[
 t(C_m,\widetilde W_\varepsilon)\leq G_m(p_k)+2m\varepsilon^2,
\]
so $\widetilde W_\varepsilon$ satisfies \eqref{eq:hypotheses} with
$\delta_e=0$ and $\delta_C=2m\varepsilon^2$. By the triangle inequality,
\Cref{lem:distance} and $\norm\cdot_\square\leq\norm\cdot_1$,
\[
 \dsq(\widetilde W_\varepsilon,W_k)
 \geq(q+\varepsilon)\varepsilon-2\varepsilon^2
 =q\varepsilon-\varepsilon^2\geq\frac q2\varepsilon
 =\frac{q}{2\sqrt{2m}}\sqrt{\delta_C},
\]
using $\varepsilon\leq\varepsilon_0\leq q/2$. Again no exponent
$c>1/2$ can hold.

Since both coordinate axes already defeat every $c>1/2$, so does the
joint statement.
\end{proof}

\begin{remark}
The two families are complete $k$-partite graphons, hence at distance
zero from the family of \emph{all} complete $k$-partite graphons with
unrestricted part measures; the distance of order $\varepsilon$ in
\Cref{lem:distance} is a distance to its balanced member only. The square
root in \Cref{thm:main} is therefore attributable to a single degenerate
direction, namely the part sizes, along which the functional
$t(C_m,\cdot)$ is quadratically rather than linearly sensitive. No
linear estimate to the unrestricted family is asserted here.
\end{remark}

\bibliographystyle{abbrv}
\bibliography{references}

\end{document}
